\Article{358.\ 359}
\textsc{Theorema.} \emph{Sit $P$ functio prima et $x^\nu$ infima potestas ipsius $x$, quae per $P$ divisa dat residuum $1$; porro sit $P\equiv(P,\rho^n)$, erit $n$ alicui numeri $p$ potestati secundum $\nu$ congruus.}

\emph{Dem.} Supra ostendimus, si $P$ sit
\[
=x^m+Ax^{m-1}+Bx^{m-2}+\text{etc.}
\]
fore
\[
\begin{aligned}
z^m+Az^{m-1}+Bz^{m-2}+\text{etc.}\qquad\qquad\qquad\qquad&\\
{}-(z-x)(z-x^p)\dotstwo(z-x^{p^{m-1}})&
\end{aligned}
\]
per $P$ divisibilem. Simili modo sequeretur esse
\[
\begin{aligned}
z^m+Az^{m-1}+Bz^{m-3}+\text{etc.}\qquad\qquad\qquad\qquad&\\
{}-(z-x^n)(z-x^{np})\dotstwo(z-x^{np^{m-1}})&
\end{aligned}
\]
per $P$ divisibilem. Quoniam autem hi factores inter se sunt primi, necessario singuli singulis secundum $P,p$ congrui esse debent. Quare $z-x^n$ debet esse $\equiv z-x^{p^\kappa}$ i. e. $p^\kappa\equiv n(\mathrm{mod.}\,\nu)$. Q. E. D.\begingroup\renewcommand{\thefootnote}{*)}\footnote{Si $(P,\rho^a)\equiv(P,\rho^b)(\mathrm{mod.}\,p)$ erit $a\equiv p^\kappa b(\mathrm{mod.}\,\nu)$.

\emph{Demonstratio.} Sit $z^m+Az^{m-1}+Bz^{m-2}+\dotstwo=\Pi$ erit $(\Pi,\rho^a)\equiv(\Pi,\rho^b)(\mathrm{mod.}\,P)$; est autem
\[
\begin{aligned}
(\Pi,\rho^a)&\equiv(z-x^a)(z-x^{ap})(z-x^{app})\dotstwo(z-x^{ap^{m-1}}),\\
(\Pi,\rho^b)&\equiv(z-x^b)(z-x^{bp})(z-x^{bpp})\dotstwo(z-x^{bpm-1})
\end{aligned}
\]
unde patet propositio.

Productum ex $\Pi,(\Pi,\rho^2),(\Pi,\rho^3)$ etc. $(\Pi,\rho^\nu)$ est $\equiv(z^\nu-1)^m(\mathrm{mod.}\,P)$; est enim
\[
\begin{aligned}
(z-x)(z-x^2)(z-x^3)\ldots(z-x^\nu)&\\
{}\equiv(z-x^p)(z-x^{2p})(z-x^{3p})\dotstwo(z-x^{\nu p})&\equiv\text{etc.}\equiv z^\nu-1
\end{aligned}
\]
In serie $P,(P,\rho^2),(P,\rho^3)$ etc.$\dotsfour(P,\rho^\nu)$ omnes divisores primi functionis $x^\nu-1$ occurrunt, et quidem quisque $m$ vicibus. Inde patet, productum ex omnibus esse $\equiv(x^\nu-1)^m$.}\endgroup

\par\medskip\Needspace{5\baselineskip}
\begin{center}\emph{De inventione divisorum primorum functionis $x^\nu-1$ secundum modulum primum.}\end{center}
\Article{360}
Si $\nu$ per modulum $p$ seu per aliquam eius potestatem est divisibilis, sit $\nu=p^k\lambda$, eritque
\[
x^\nu-1\equiv(x^\lambda-1)^{p^k}(\mathrm{mod.}\,p)
\]
Unde manifestum est, eum tantummodo casum considerari oportere, ubi $\nu$ per $p$ non dividitur.

\SourceLeaf{240}
Si $p^m\equiv1(\mathrm{mod.}\,\nu)$ et quidem $m$ quam minimus, tum patet $x^{p^m-1}-1$ per $x^\nu-1$ dividi posse. Quamobrem $x^\nu-1$ alios divisores habere nequit quam $x^{p^m-1}-1$. At haecce expressio habet divisores primos $m$ dimensionum aliosque, quorum dimensionum numerus est divisor numeri $m$. Tales igitur etiam $x^\nu-1$ habebit. Quot autem cuiusvis generis habeat, per exemplum declaramus, unde facile lex generalis deduci poterit.

Sit $\nu=63$ et $p=13$, erit $m=6$. Quare $x^{63}-1$ secundum modulum $13$ factores primos habebit sex, trium, duarum dimensionum uniusque. Iam palam est, productum ex factoribus unius dimensionis fore divisorem communem (maximae dimensionis) functionum $x^{63}-1$ et $x^{12}-1$ i. e. $x^3-1$; quare tres erunt factores primi unius dimensionis. Productum ex omnibus factoribus primis duarum dimensionum uniusque erit divisor communis functionum $x^{63}-1$ et $x^{168}-1$ i. e. $x^{21}-1$; quare erunt $\frac{21-3}2$ sive $9$ factores duarum dimensionum. Productum ex factoribus primis trium dimensionum uniusque erit divisor communis functionum $x^{63}-1$ et $x^{2196}-1$ i. e. $x^9-1$; quare erunt $\frac{9-3}3$ i. e. $2$ divisores trium dimensionum. Tandem reliqui erunt sex dimensionum, quorum igitur numerus $=\frac{63-6-18-3}6$ i. e. $6$.

Facile per attentam huius rei ponderationem sequens regula generalis deducitur:

\emph{Sit $\delta$ divisor ipsius $m$, sint omnes numeri $\delta$ divisores ipso $\delta$ minores $\delta',\delta'',\delta'''$ etc. Sint divisores communes maximi ipsius $\nu$ cum $p^\delta-1$, $p^{\delta'}-1$, $p^{\delta''}-1$ etc. respective $\mu,\mu',\mu''$ etc., sit $\frac\mu{\mu'};\frac\mu{\mu''},\frac\mu{\mu'''}$ etc. $=\lambda',\lambda'',\lambda'''$ etc. habebitque $x^\nu-1$ $\frac1\delta$ties tot divisores primos $\delta$ dimensionum, quot infra numerum $\mu$ sunt numeri per nullum numerorum $\lambda',\lambda'',\lambda'''$ etc. divisibiles.}

\Article{361}
\textsc{Theorema.} \emph{Si functio $X$ indeterminatae $x$ per aliam $\xi$ dividi possit et $X$ si pro $x$ scribatur $x^k$, transeat in $X'$, $X'$ per $(\xi,\rho^{\frac1k})$ dividi poterit.}

\emph{Dem.} Sit $X\equiv\xi\upsilon$ transeantque $\xi,\upsilon$ in $\xi',\upsilon'$, si pro $x$ scribatur $x^k$. Patet, fore $X'\equiv\xi'\upsilon'$. At $\xi'$ per $(\xi,\rho^{\frac1k})$ dividi potest. Quare etiam $X'$. Q. E. D.
